\documentclass[12pt]{article}

\usepackage[utf8]{inputenc}
\usepackage{amsmath, amssymb, amsthm}
\usepackage{geometry}
\geometry{a4paper, margin=2.5cm}

\title{Existența Primitivei pentru Funcții Olomorfe}
\author{}
\date{}

\theoremstyle{definition}
\newtheorem{definition}{Definiție}

\theoremstyle{plain}
\newtheorem{theorem}{Teoremă}
\newtheorem{proposition}{Propoziție}
\newtheorem{corollary}{Corolar}

\begin{document}

\maketitle

\begin{definition}
Fie $\Omega \subset \mathbb{C}$ un domeniu.  
O funcție $F:\Omega \to \mathbb{C}$ se numește \emph{primitivă} a funcției $f:\Omega \to \mathbb{C}$ dacă


\[
F'(z) = f(z) \quad \text{pentru toate } z \in \Omega.
\]


\end{definition}

\begin{theorem}
(\textbf{Teorema lui Cauchy privind existența primitivei})  
Fie $D = D(z_0, R)$ un disc deschis în $\mathbb{C}$ și fie $f : D \to \mathbb{C}$ o funcție olomorfă.  
Atunci $f$ admite o primitivă în $D$.  
Mai precis, funcția


\[
F(z) = \int_{\gamma_{z_0,z}} f(w)\,dw,
\]


unde $\gamma_{z_0,z}$ este orice curbă netedă din $z_0$ la $z$ în interiorul discului, este bine definită și satisface


\[
F'(z) = f(z).
\]


\end{theorem}

\begin{proposition}
(\textbf{Unicitatea primitivei})  
Dacă $F$ și $G$ sunt două primitive ale aceleiași funcții olomorfe $f$ într-un disc, atunci


\[
F(z) - G(z) = \text{constantă}.
\]


\end{proposition}

\begin{corollary}
(\textbf{Integrala unei funcții olomorfe pe o curbă închisă dintr-un disc})  
Fie $D = D(z_0, R)$ un disc și $f : D \to \mathbb{C}$ o funcție olomorfă.  
Atunci, pentru orice curbă închisă netedă $\gamma : [a,b] \to D$, are loc


\[
\int_{\gamma} f(z)\,dz = 0.
\]


\textit{În special, orice funcție olomorfă pe un disc are integrală nulă pe orice contur închis conținut în disc.}
\end{corollary}

\end{document}
